\documentclass{article}
\usepackage[utf8]{inputenc}
\usepackage[T1]{fontenc}
\usepackage{amsmath}
\usepackage{amsfonts}
\usepackage{geometry}
\geometry{margin=2.5cm}
\begin{document}

\title{GlobalFlow Network Design Optimization Model}
\author{Phase 1: Baseline Scenario}
\date{}
\maketitle

\section{Ensembles}

\begin{itemize}
    \item $p \in P = \{A, B, C\}$ : les produits (A = Fertilizers, B = Semiconductors, C = Battery Components)
    \item $s \in S = \{S1, \dots, S9\}$ : l'ensemble des fournisseurs (suppliers)
    \item $w \in W = \{W1, \dots, W15\}$ : l'ensemble des entrepôts (regional warehouses)
    \item $h \in H = \{H1, \dots, H5\}$ : l'ensemble des hubs (Frankfurt, Dubai, Singapore, Chicago, São Paulo)
    \item $c \in C = \{C1, \dots, C30\}$ : l'ensemble des clients (customers)
    \item $n \in N = \{1, \dots, 59\}$ : l'ensemble des nœuds ($N = S \cup W \cup H \cup C$)
    \item $a \in A = \{A1, \dots, A317\}$ : l'ensemble des arcs du réseau
\end{itemize}

\subsection{Sous-ensembles d'arcs}

\begin{itemize}
    \item $a \in A_{\text{fixed}} \subset A$ : l'ensemble des arcs avec coût d'activation > 0 (47 arcs)
    \item $a \in A_{\text{always}} \subset A$ : l'ensemble des arcs sans coût d'activation (270 arcs).
    \quad Remarque: $A_{\text{fixed}} \cup A_{\text{always}} = A$ et $A_{\text{fixed}} \cap A_{\text{always}} = \emptyset$
\end{itemize}

\subsection{Sous-ensembles de fournisseurs par produit}

Chaque fournisseur ne produit qu'un seul type de produit. On définit donc :

\begin{itemize}
    \item $S^p \subset S$ : l'ensemble des fournisseurs qui produisent le produit $p \in P$.
    \quad Formellement : $S^p = \{ s \in S \mid Sup_s^p > 0 \}$.
    \item Les ensembles $S^p$ forment une partition de $S$ :
    $\displaystyle\bigcup_{p \in P} S^p = S$ \quad et \quad $S^p \cap S^{p'} = \emptyset \; \forall p \neq p'$.
\end{itemize}

\noindent Dans les données :
$S^A = \{S1, S2, S3\}$,\quad
$S^B = \{S4, S5, S6\}$,\quad
$S^C = \{S7, S8, S9\}$.

\subsection{Arcs par nœud ou type de nœud}

\begin{itemize}
    \item $a \in S_{\text{source}}^n \subset A$ : ensemble des arcs $a$ dont la source est le nœud $n \in N$.
    \item $a \in S_{\text{target}}^n \subset A$ : ensemble des arcs $a$ dont la destination est le nœud $n \in N$.
    \item $a \in S_{\text{source}}^s \subset A$ : ensemble des arcs $a$ dont la source est un fournisseur $s \in S$.
    \item $a \in S_{\text{target}}^c \subset A$ : ensemble des arcs $a$ dont la destination est un client $c \in C$.
    \item $a \in S_{\text{source}}^w \subset A$ : ensemble des arcs $a$ dont la source est un entrepôt $w \in W$.
    \item $a \in S_{\text{target}}^w \subset A$ : ensemble des arcs $a$ dont la destination est un entrepôt $w \in W$.
    \item $a \in S_{\text{source}}^h \subset A$ : ensemble des arcs $a$ dont la source est un hub $h \in H$.
    \item $a \in S_{\text{target}}^h \subset A$ : ensemble des arcs $a$ dont la destination est un hub $h \in H$.
\end{itemize}

\subsection{Structure de l'arc}

Chaque arc $a \in A$ est associé avec :
\begin{itemize}
    \item $i(a) \in N$ : l'origine (source) de l'arc.
    \item $j(a) \in N$ : la destination (target) de l'arc.
\end{itemize}

\section{Paramètres}

\begin{itemize}
    \item $Dem_c^p$ : Demande du produit $p$ au client $c$ (unités).
    \item $Sup_s^p$ : Production disponible du produit $p$ au fournisseur $s$ (unités).
    \item $OpenCost_w$ : Coût d'ouverture de l'entrepôt $w$ (monétaire).
    \item $CapWare_w$ : Capacité de l'entrepôt $w$ (unités).
    \item $CapArc_a$ : Capacité partagée sur l'arc $a$ pour tous les produits (unités).
    \item $FixActCost_a$ : Coût fixe d'activation de l'arc $a$
    \quad (monétaire, = 0 pour $a \in A_{\text{always}}$, > 0 pour $a \in A_{\text{fixed}}$).
    \item $VarArcCost_a^p$ : Coût variable de passage du produit $p$ sur l'arc $a$ (monétaire par unité, avant tarif).
    \item $zone\_from_a$ : la zone tarifaire d'origine de l'arc $a$.
    \item $zone\_to_a$ : la zone tarifaire de destination de l'arc $a$.
    \item $TariffRate_{z_1, z_2}$ : Taux de tarif douanier pour les arcs reliant la zone $z_1$ à la zone $z_2$.
    \quad Déterminé en cherchant la paire $(zone\_from_a, zone\_to_a)$ dans la table TariffZones;
    \quad par défaut = 0 si la paire n'existe pas.
    \item $TariffRate_a$ : Le taux de tarif applicable à l'arc $a$, déterminé par la paire
    \quad $(zone\_from_a, zone\_to_a)$.
    \item $TotalCost_a^p$ : Coût total de passage du produit $p$ sur l'arc $a$, incluant le tarif douanier.

    \[
        TotalCost_a^p = (1 + TariffRate_a) \cdot VarArcCost_a^p \quad \forall a \in A, p \in P
    \]

    où $TariffRate_a = TariffRate_{(zone\_from_a, zone\_to_a)}$ (lookup dans la table TariffZones, défaut 0).

\end{itemize}

\section{Variables de décision}

\begin{itemize}
    \item $x_a^p \ge 0$ : Quantité de produit $p$ passant par l'arc $a$ (unités).
    La variable $x_a^p$ n'est définie que si l'arc $a$ n'est pas un arc fournisseur,
    ou si $a$ est un arc fournisseur avec $i(a) \in S^p$ (le fournisseur produit $p$).
    Formellement : $x_a^p = 0$ est imposé implicitement lorsque $i(a) \in S$ et $i(a) \notin S^p$.

    \item $openWarehouse_w \in \{0, 1\}$ : Variable binaire indiquant si l'entrepôt $w$ est ouvert.
    \begin{equation*}
        openWarehouse_w = \begin{cases}
            1 & \text{si l'entrepôt } w \text{ est ouvert} \\
            0 & \text{sinon}
        \end{cases}
    \end{equation*}

    \item $arc_a \in \{0, 1\}$ : Variable binaire indiquant si l'arc $a$ est activé.
    \quad Défini uniquement pour $a \in A_{\text{fixed}}$.
    \begin{equation*}
        arc_a = \begin{cases}
            1 & \text{si l'arc } a \text{ est activé} \\
            0 & \text{sinon}
        \end{cases} \quad \forall a \in A_{\text{fixed}}
    \end{equation*}

\end{itemize}

\section{Contraintes}

\subsection{Satisfaction de la demande et disponibilité de la production}

\noindent\textbf{(C1) Satisfaction de la demande :}
Pour chaque produit $p$ et chaque client $c$, la somme des flux entrants doit égaler la demande.
\begin{equation}
    \sum_{a \in S_{\text{target}}^c} x_a^p = Dem_c^p \quad \forall p \in P, \forall c \in C \label{eq:demand}
\end{equation}

\noindent\textbf{(C2) Contrainte de production :}
La somme des flux sortants d'un fournisseur $s$ pour le produit $p$ qu'il produit ne peut dépasser
sa production disponible. La contrainte est définie uniquement pour les paires $(s, p)$ telles que
$s \in S^p$ (le fournisseur produit ce produit).
\begin{equation}
    \sum_{a \in S_{\text{source}}^s} x_a^p \le Sup_s^p \quad \forall p \in P,\; \forall s \in S^p
    \label{eq:supply}
\end{equation}

\noindent\textit{Remarque :} L'inégalité ($\le$) au lieu de l'égalité ($=$) est justifiée par
deux raisons : (i) les données présentent de légères asymétries entre offre totale et demande totale
(e.g.\ $\sum_{s \in S^B} Sup_s^B = 2223 > 2222 = \sum_c Dem_c^B$), et (ii) la combinaison de
la contrainte (C1) d'égalité sur la demande et des contraintes de conservation du flux (C6)-(C7)
garantit que tout flux sortant d'un fournisseur atteint effectivement un client.
De plus, la restriction $s \in S^p$ (via les sous-ensembles $S^p$) empêche tout flux
fantôme de produits qu'un fournisseur ne produit pas.

\subsection{Capacités et activation des arcs}

\noindent\textbf{(C3) Capacité des arcs toujours actifs :}
Les arcs sans coût d'activation ont une capacité fixe.
\begin{equation}
    \sum_{p \in P} x_a^p \le CapArc_a \quad \forall a \in A_{\text{always}} \label{eq:arc_cap_always}
\end{equation}

\noindent\textbf{(C4) Capacité des arcs optionnels liée à l'activation :}
Les arcs avec coût d'activation ne peuvent transporter du flux que s'ils sont activés.
\begin{equation}
    \sum_{p \in P} x_a^p \le CapArc_a \cdot arc_a \quad \forall a \in A_{\text{fixed}} \label{eq:arc_cap_fixed}
\end{equation}

\subsection{Capacités et activation des entrepôts}

\noindent\textbf{(C5) Capacité de l'entrepôt liée à l'ouverture :}
L'inflow total vers un entrepôt $w$ ne peut excéder sa capacité, et seulement si l'entrepôt est ouvert.
\begin{equation}
    \sum_{a \in S_{\text{target}}^w} \sum_{p \in P} x_a^p \le CapWare_w \cdot openWarehouse_w \quad \forall w \in W \label{eq:wh_cap}
\end{equation}

\subsection{Conservation du flux}

\noindent\textbf{(C6) Conservation du flux aux entrepôts :}
Tout produit entrant dans un entrepôt doit en ressortir.
\begin{equation}
    \sum_{a \in S_{\text{source}}^w} x_a^p = \sum_{a \in S_{\text{target}}^w} x_a^p \quad \forall p \in P, \forall w \in W \label{eq:wh_flow}
\end{equation}

\noindent\textbf{(C7) Conservation du flux aux hubs :}
Tout produit entrant dans un hub doit en ressortir (les hubs sont des points de transshipment passants).
\begin{equation}
    \sum_{a \in S_{\text{source}}^h} x_a^p = \sum_{a \in S_{\text{target}}^h} x_a^p \quad \forall p \in P, \forall h \in H \label{eq:hub_flow}
\end{equation}

\subsection{Domaine des variables}

\noindent\textbf{(C8) Non-négativité :}
\begin{align}
    x_a^p &\ge 0 \quad \forall a \in A, \forall p \in P \label{eq:flow_nonneg}
\end{align}

\noindent\textbf{(C9) Intégrité et domaine :}
\begin{align}
    openWarehouse_w &\in \{0, 1\} \quad \forall w \in W \label{eq:wh_binary}\\
    arc_a &\in \{0, 1\} \quad \forall a \in A_{\text{fixed}} \label{eq:arc_binary}
\end{align}

\section{Fonction objectif}

Minimiser le coût total du réseau, comprenant les coûts fixes d'ouverture d'entrepôts,
les coûts fixes d'activation d'arcs, et les coûts variables de transport (incluant les tarifs) :

\begin{equation}
    \min \quad
    \sum_{w \in W} OpenCost_w \cdot openWarehouse_w
    + \sum_{a \in A_{\text{fixed}}} FixActCost_a \cdot arc_a
    + \sum_{a \in A} \sum_{p \in P} TotalCost_a^p \cdot x_a^p
    \label{eq:objective}
\end{equation}

\noindent où $TotalCost_a^p = (1 + TariffRate_a) \cdot VarArcCost_a^p$.

\section{Remarques de modélisation}

\begin{itemize}
    \item \textbf{Partition des arcs :} L'ensemble $A$ est partitionné en $A_{\text{fixed}}$ (47 arcs avec coûts
    d'activation) et $A_{\text{always}}$ (270 arcs toujours actifs). Les contraintes de capacité et les variables
    binaires d'activation $arc_a$ ne s'appliquent que sur $A_{\text{fixed}}$.

    \item \textbf{Tarifs douaniers :} Le taux $TariffRate_a$ est déterminé par la paire de zones
    $(zone\_from_a, zone\_to_a)$ en cherchant dans la table TariffZones. Si la paire n'existe pas,
    le taux par défaut est 0 (pas de tarif).

    \item \textbf{Coûts de transport :} Le coût total $TotalCost_a^p$ intègre le coût variable base et
    la majoration tarifaire selon la formule $(1 + TariffRate_a) \cdot VarArcCost_a^p$.

    \item \textbf{Hubs non capacités :} Les hubs respectent la conservation du flux mais ne sont pas
    soumis à une contrainte de capacité (points de transshipment de capacité infinie).

    \item \textbf{Entrepôts capacités et activables :} Les entrepôts régionaux sont activables via
    la variable binaire $openWarehouse_w$. Un entrepôt ne peut recevoir du flux que s'il est ouvert,
    et l'inflow total est limité par sa capacité.

    \item \textbf{Sous-ensembles de fournisseurs par produit ($S^p$) :}
    Chaque fournisseur ne produisant qu'un seul type de produit, on définit $S^p \subset S$
    comme l'ensemble des fournisseurs de $p$. La contrainte (C2) est restreinte à $s \in S^p$
    et les variables de flux $x_a^p$ ne sont pas créées pour les arcs fournisseurs incompatibles
    ($i(a) \in S \setminus S^p$). Ceci évite les flux fantômes : sans cette restriction,
    le solveur pourrait faire transiter du produit $p$ depuis un fournisseur qui ne le produit pas,
    contournant ainsi les contraintes d'offre.

    \item \textbf{Contrainte de production en inégalité :}
    La contrainte (C2) utilise $\le$ plutôt que $=$. En présence d'un léger déséquilibre
    offre/demande ($\sum_{s \in S^p} Sup_s^p \ge \sum_c Dem_c^p$), l'égalité rendrait
    le modèle infaisable. La combinaison de (C1) (demande exacte), (C6)-(C7) (conservation)
    et (C2) ($\le$) garantit que la production est utilisée au juste niveau pour satisfaire
    la demande, sans surplus bloqué dans le réseau.

\end{itemize}

\end{document}